\subsection{将 G 的线性表示延拓到 G 上的测度}

设 G 为局部紧群，E 为局部凸空间，U 为 G 在 E 中的线性表示。假设 U 连续且 E 拟完备。则对每个测度 \(\mu \in \mathcal{C}'(G)\)，有

\[
\int_ {\mathbf {G}} U (s) d \mu (s) \in \mathcal {L} (\mathbf {E}; \mathbf {E})
\]

（Ch. VI, §1, No.~7）。记 \(U(\mu) = \int_{\mathbf{G}} U(s) d\mu(s)\)。在 \(\mathcal{C}'(\mathbf{G})\) 上赋予 \(\mathcal{C}(\mathbf{G})\) 中的紧收敛拓扑。映射 \((\mu, x) \mapsto U(\mu)x\) 从 \(\mathcal{C}'(\mathbf{G}) \times \mathbf{E}\) 到 \(\mathbf{E}\) 相对于 \(\mathcal{C}'(\mathbf{G})\) 的等度连续子集和 \(\mathbf{E}\) 的紧子集是拟连续的；特别地，映射 \(\mu \mapsto U(\mu)\) 从 \(\mathcal{C}'(\mathbf{G})\) 到 \(\mathcal{L}(\mathbf{E}; \mathbf{E})\)（在后者赋予紧收敛拓扑）是连续的（同上，命题 16）。

为了以后能够应用这些结果，我们注意到，若 X 是局部紧空间，则赋予紧收敛拓扑的 \(\mathcal{C}(\mathrm{X})\) 是完备的（GT, X, §1, No.~6, Th. 2 的推论 3）。另一方面，\(\mathcal{K}(\mathrm{X})\) 是桶状的，因而其对偶 \(\mathcal{M}(\mathrm{X})\) 在 \(\mathcal{K}(\mathrm{X})\) 中赋予紧收敛拓扑后是拟完备的（TVS, III, §4, No.~2, Th. 1 的推论 4）。当然，\(\overline{\mathcal{K}(\mathrm{X})}\) 关于由其范数导出的拓扑是完备的，因而其对偶 \(\mathcal{M}^{1}(\mathrm{X})\) 关于 \(\overline{\mathcal{K}(\mathrm{X})}\) 中的紧收敛拓扑是拟完备的（同上）。

现在假设 U 是局部紧群 G 在 Banach 空间 E 上的连续线性表示。置 \(g(s) = \|U(s)\|\)，对所有 \(s \in G\) 而言。于是，若 \(\mu\) 是 G 上的测度且 g 是 \(\mu\)-可积的，则 \(\int_{\mathbf{G}} U(s) d\mu(s) \in \mathcal{L}(\mathbf{E}; \mathbf{E})\)，并且 \(\left\|\int_{\mathbf{G}} U(s) d\mu(s)\right\| \leqslant \int g(s) d|\mu|(s)\)（Ch. VI, §1, No.~7, 注 1）。仍记 \(U(\mu) = \int_{\mathbf{G}} U(s) d\mu(s)\)。

\begin{enumerate}
\def\labelenumi{\arabic{enumi}.}
\setcounter{enumi}{6}
\tightlist
\item
  自同态 \(U(\mu)\) 与自同态 \(U(s)\) 之间的关系
\end{enumerate}

引理 4. --- 设 T 为局部紧空间，a 为 T 中的一点，M 为 \(\mathcal{M}(\mathrm{T})\) 的子集，\(\mathfrak{F}\) 为 M 上的滤子。假设：

\begin{enumerate}
\def\labelenumi{(\roman{enumi})}
\item
  对每个紧子集 \(\mathbf{K}\)，它是 \(\mathbf{T}\) 的子集，数 \(|\mu|(\mathbf{K})\) 在 \(\mu \in \mathbf{M}\) 时有上界；
\item
  \(\lim_{\mu, \mathfrak{F}} |\mu| (\mathrm{K}) = 0\) 对 \(\mathrm{K}\) 为 \(\mathrm{T} - \{a\}\) 的每个紧子集成立。
\item
  存在 T 中 a 的紧邻域 V，使得 \(\lim_{\mu, \mathfrak{F}} \mu(V) = 1\)。
\end{enumerate}

则滤子 \(\mathfrak{F}\) 收敛于 \(\varepsilon_{a}\)；且 \(\mathcal{M}(\mathrm{T})\) 赋有 \(\mathcal{H}(\mathrm{T})\) 中的紧收敛拓扑。

由假设 (i)，M 是 \(\mathcal{M}(\mathrm{T})\) 的等度连续子集，因为它是弱有界的，而 \(\mathcal{H}(\mathrm{T})\) 是桶状的（TVS, III, §4, No.~2, Th. 1）。因此，只需（GT, X, §2, No.~4, Th. 1）证明：若 \(f \in \mathcal{H}(\mathrm{T})\) ，则 \(\lim_{\mu, \mathfrak{F}} \mu(f) = f(a)\) 。令 K 为 V 与 \(f\) 的支集之并；若 \(\mathrm{K}'\) 是 \(\mathrm{K} - \mathrm{V}\) 的闭包，则有

\[
| \mu (\mathrm{K}) - \mu (\mathrm{V}) | = | \mu (\mathrm{K} - \mathrm{V}) | \leqslant | \mu | (\mathrm{K} ^ {\prime});
\]

由于 \(\mathbf{K}'\) 是紧的且不包含 \(a\) ，由此可得 \(\lim_{\mu, \mathfrak{F}} \mu(\mathbf{K}) = 1\) 。令 \(\varepsilon > 0\) ，并令 \(\mathbf{W}\) 为 \(a\) 的开邻域，位于 \(\mathbf{K}\) 中，使得 \(|f(t) - f(a)| \leqslant \varepsilon\) 对 \(t \in \mathbf{W}\) 成立；可写成

\[
\mu (f) - f (a) = f (a) \big (\mu (\mathrm{K}) - 1 \big) + \int_ {\mathrm{K}} \big (f (t) - f (a) \big) d \mu (t);
\]

关于 K 的积分可写为 W 与 K - W 上相应积分之和；若 \(C = \sup |f|\) ，则有

\[
| \mu (f) - f (a) | \leqslant \mathrm{C} | \mu (\mathrm{K}) - 1 | + \varepsilon \cdot | \mu | (\mathrm{K}) + 2 \mathrm{C} \cdot | \mu | (\mathrm{K} - \mathrm{W}).
\]

由于右端第一项和第三项相对于 \(\mathfrak{F}\) 趋于 0，故确实有 \(\lim_{\mu \to \mathfrak{F}} \mu(f) = f(a)\) 。

推论 1. --- 在引理 4 的假设下，再假定存在 T 的一个紧子集 \(K_{0}\)，包含所有测度 \(\mu \in M\) 的支集。则 F 也收敛于 \(\varepsilon_{a}\)；在 \(\mathcal{C}'(T)\) 上赋予 \(\mathcal{C}(T)\) 中的紧收敛拓扑。

事实上，从 \(\mathcal{C}(\mathrm{T})\) 到 \(\mathcal{C}(\mathrm{K}_{0})\) 的限制映射是连续的；因此，若 H 是 \(\mathcal{C}(\mathrm{T})\) 的紧子集，则 H 中各函数限制在 \(K_{0}\) 上后构成 \(\mathcal{C}(\mathrm{K}_{0})\) 的紧子集。于是只需在以 \(K_{0}\) 替代 T 后应用引理 4。

推论 2. --- 在推论 1 的假设下，设 f 是从 T 到拟完备局部凸空间 E 的连续映射。则

\[
\lim _ {\mu , \mathfrak {F}} \int f (t) d \mu (t) = f (a).
\]

这由推论 1 以及第 VI 章 §1 No.~6 的命题 14 得出。

推论 3. --- 设 G 是局部紧群，E 是拟完备局部凸空间，U 是 G 在 E 中的连续线性表示。设 \(\beta\) 是 G 上的正测度，a 是 G 的一个元素，B 是 a 的邻域滤子的一个基，由紧邻域组成。对每个 \(V \in B\)，令 \(f_{V}\) 是 G 上满足 \(\geqslant 0\) 的连续函数，其支集包含于 V 中，并且满足 \(\int f_{V} d\beta = 1\)。则对每个 \(x \in E\)，

\[
U (a) x = \lim _ {\mathrm{V}} U (f _ {\mathrm{V}} \cdot \beta) x,
\]

这里的极限是相对于 \(\mathfrak{B}\) 的截面滤子取得的。

映射 \(s \mapsto U(s)x\) 从 \(G\) 到 \(E\) 是连续的。由推论 2，\(U(a)x = \lim_{V} \int (U(s)x) \cdot f_V(s) d\beta(s)\)，其中极限是相对于 \(\mathfrak{B}\) 的截面滤子取得的；也就是说，\(U(a)x = \lim_{V} U(f_V \cdot \beta)x\)。

命题 10. --- 设 G 是局部紧群，E 是拟完备局部凸空间，U 是 G 在 E 中的连续线性表示，且 \(\beta\) 是支集为 G 的 G 上正测度。

\begin{enumerate}
\def\labelenumi{(\roman{enumi})}
\item
  向量 \(U(f \cdot \beta)x\)，其中 \(f\) 遍历 \(\mathcal{K}(\mathbf{G})\) 且 \(x\) 遍历 \(\mathbf{E}\)，在 \(\mathbf{E}\) 中稠密。
\item
  设 \(\mathbf{F}\) 是 \(\mathbf{E}\) 的闭线性子空间。若 \(\mathbf{F}\) 对 \(U\) 稳定，则 \(U(\mu)(\mathbf{F}) \subset \mathbf{F}\) 对每个 \(\mu \in \mathcal{C}'(\mathbf{G})\) 成立。反之，若 \(U(f \cdot \beta) \subset \mathbf{F}\) 对每个 \(f \in \mathcal{K}(\mathbf{G})\) 成立，则 \(\mathbf{F}\) 对 \(U\) 稳定。
\end{enumerate}

（ii）的第一部分是显然的，因为将 \(U(s)\) 限制在 F 上，对于 \(s \in G\)，便定义了 G 在拟完备局部凸空间 F 中的连续线性表示。（ii）的第二部分和（i）由引理 4 的推论 3 得出。

